\scanpage{113}
$\tau_\infty(H^{**}(X))$ [qui est une algèbre graduée \deleted{\unsure{non}} alternée, entre autres, de $M_\infty$] est \deleted{\unsure{provient donc}} dans l'image essentielle de $\mathrm{Mod}(k)$, i.e. peut s'écrire (de façon \unsure{can.} \unsure{une fois qu'on en a un modèle})
\[
\mathcal H_C^{**}(X)=H_C^*(X)\otimes_k S_k
\]
\unsure{Et on sait} que $H_C^*(X)$ hérite de la structure de $k$-algèbre graduée de $\tau_\infty(H^{**}(X))$. Cela implique que \deleted{pour} les $C_n^{**}$ ont la $=$ propriété (faire $X=\Delta_n$).
\deleted{On veut de plus que \inserted{$\forall X$,} l'augmentation $S_k\longrightarrow C^{**}(X)$ \unsure{définisse}}
\unsure{On va préciser} cette exigence.

On peut considérer $\unsure{C_n^{**}|C_*^{**}}$ comme un faisceau (à structures de type $\mathcal D$) sur le topos $(\mathrm{Ss})$, il est coaugmenté par le faisceau constant
\[
S_k\longrightarrow C_*^{**}
\]
\uline{[NB}, on regarde ici le degré extérieur comme degré principal] de sorte qu'on a un diagramme (dans la cat. dérivée \deleted{\unsure{convenable d'abord}} \inserted{$\unsure{D(S_k\text{-mod grad})}$}) :
\[
(1)\qquad
\diagram{\begin{tikzcd}[column sep=large]
\mathrm R^*\Gamma(X,S_k)\arrow[r]&\mathrm R^*\Gamma(X,C_*^{**})\\
& C^{**}(X)=\Gamma(X,C_*^{**}\arrow[u]
\end{tikzcd}}
\]
\margin{112}
